\chapter{Robustness Bubbles}
\label{ch:robustness}

This chapter examines whether the bubble results in \cref{ch:results} are sensitive to three implementation choices in the PSY/GSADF design. \Cref{sec:rob-lags} varies the fixed lag order in the recursive ADF regressions. \Cref{sec:rob-wild} compares the corrected wild-multiplier bootstrap used in the main analysis with the multiplier implementation shipped in \texttt{psymonitor}~0.0.3. \Cref{sec:rob-log} re-estimates the test variables in logs rather than levels. The objective is not to generate a new interpretation, but to test whether the main conclusions depend on lag choice, bootstrap implementation, or the level-form specification. Throughout the chapter, the sample, recursive window rule, control window, bootstrap count, and date-stamping logic follow \cref{subsec:bubble-specification,subsec:bubble-critical-values,subsec:bubble-date-stamping}; only the robustness dimension under review is varied.


% ---------------------------------------------------------------------
\section{Lag-order robustness}
\label{sec:rob-lags}

The baseline specification in \cref{ch:results} follows the PSY convention of reporting the empirical application at $k=0$, while \citet[footnote~22]{phillipsShiYu2015a} use $k=3$ as a principal alternative. To examine whether the conclusions are driven by the no-lag specification, the analysis is repeated for $k\in\{1,2,3\}$, with wild-multiplier bootstrap critical values recomputed separately at each lag. The GSADF statistic at a given lag is therefore always compared with the corresponding finite-sample null distribution at that same lag. \Cref{tab:rob-lags-headline} reports the sixteen headline series; the concentration-basket counterpart is reported in \cref{tab:rob-lags-topn}.

The headline evidence is robust in the dimensions that matter for the thesis, but it is not mechanically invariant across every series. Twelve of the sixteen headline series reject the unit-root-with-drift null at the 90\,\% bootstrap critical value at every lag. The four exceptions are all secondary S\&P-normalised diagnostic ratios: the value-weighted Moderate~AI~/~S\&P~500 ratio fails to reject at $k=3$, the value-weighted No~AI~/~S\&P~500 ratio fails at $k\in\{2,3\}$, the equal-weighted Moderate~AI~/~S\&P~500 ratio fails at $k=1$, and the equal-weighted No~AI~/~S\&P~500 ratio fails at $k=3$. By contrast, no High-AI series drops out of the rejection region, and all absolute indices reject at every lag. The lag sensitivity is therefore concentrated in the weaker breadth diagnostics rather than in the High-AI level, High-AI relative, or concentration evidence. This is also the direction predicted by the cross-sectional design: the first verdicts to become insignificant under a more demanding lag specification are the placebo-style and middle-exposure S\&P-normalised diagnostics, not the High-AI or No-AI denominator contrasts that drive the main interpretation.

Two features of \cref{tab:rob-lags-headline} are central. First, the High~AI absolute index remains significant at the 99\,\% level at every $k\in\{0,1,2,3\}$ under both value- and equal-weighting. Moreover, the GSADF statistic rises rather than falls when lags are added: under value-weighting it increases from 4.04 at $k=0$ to 5.21 at $k=2$ and 4.98 at $k=3$, while under equal-weighting it increases from 4.25 at $k=0$ to 5.08 at both $k=2$ and $k=3$. Higher-order autoregressive correction therefore does not weaken the High-AI level rejection. Second, the High~AI~/~No~AI ratio, which is the primary cross-sectional contrast in \cref{subsec:bubble-test-variables}, remains significant under both weightings at every lag: it is classified at the 99\,\% level for $k\in\{0,1,2\}$ and at the 95\,\% level for $k=3$. The exact significance bin is therefore not fully invariant, but the rejection itself is preserved across the full lag range. The one-bin decline at $k=3$ is consistent with the finite-sample efficiency loss that can arise when the recursive ADF regression is over-specified with additional lags in a relatively short minimum window.

The concentration baskets in \cref{tab:rob-lags-topn} show a similar pattern. Eleven of the twelve ratios reject at the 90\,\% level at every lag. The only ratio that drops out is the equal-weighted Top-decile~/~High~AI series, which rejects at $k\in\{0,2\}$ but fails to reject at $k\in\{1,3\}$; this was already the weakest concentration contrast in the baseline table. The Mag-7~/~No~AI ratio remains the largest concentration signal at every lag, especially under equal-weighting, where the GSADF statistic declines from 6.72 at $k=0$ to 3.23 at $k=3$ but stays in the 99\,\% rejection bin throughout. The Mag-7~/~High~AI ratios weaken as lag length increases, but both continue to reject at every lag. Thus the concentration result is robust in its ordering: the strongest evidence remains with Mag-7 and the strongest denominator contrast remains No~AI.

Read together with \cref{sec:results-headline,sec:results-concentration}, the lag exercise supports the main interpretation with an important qualification. The cross-sectional ordering High~AI~$>$~Moderate~AI~$>$~No~AI is preserved across lag choices, and the High~AI~/~No~AI and Mag-7~/~No~AI contrasts remain significant throughout. What changes with lag order is primarily the significance classification of marginal S\&P-normalised breadth diagnostics and the weakest Top-decile~/~High~AI ratio. The lag robustness exercise therefore supports the substantive results, but not the stronger claim that every reported rejection is insensitive to lag choice.

\begin{table}[t]
\centering
\caption{Lag-order robustness: headline series.}
\label{tab:rob-lags-headline}
\small
\begin{tabular}{l c S[table-format=1.3] c S[table-format=1.3] c S[table-format=1.3] c S[table-format=1.3] c c}
\toprule
 & & \multicolumn{2}{c}{$k=0$} & \multicolumn{2}{c}{$k=1$} & \multicolumn{2}{c}{$k=2$} & \multicolumn{2}{c}{$k=3$} & \\
\cmidrule(lr){3-4}\cmidrule(lr){5-6}\cmidrule(lr){7-8}\cmidrule(lr){9-10}
Series & Wgt. & {GSADF} & sig & {GSADF} & sig & {GSADF} & sig & {GSADF} & sig & all stable \\
\midrule
High AI (abs)              & vw & 4.037 & *** & 4.450 & *** & 5.214 & *** & 4.984 & *** & yes \\
Moderate AI (abs)          & vw & 2.345 & **  & 2.455 & **  & 2.801 & *** & 2.532 & **  & no  \\
No AI (abs)                & vw & 1.860 & **  & 2.057 & **  & 2.520 & **  & 2.482 & **  & yes \\
S\&P~500 (abs)             & vw & 2.167 & **  & 2.522 & *** & 2.961 & *** & 3.105 & *** & no  \\
High AI / S\&P~500         & vw & 3.408 & *** & 3.310 & *** & 3.850 & *** & 3.333 & *** & yes \\
Moderate AI / S\&P~500     & vw & 1.382 & **  & 1.357 & *   & 2.171 & **  & 1.310 & ---  & no  \\
No AI / S\&P~500           & vw & 1.453 & *   & 1.360 & *   & 1.206 & --- & 1.056 & --- & no  \\
High AI / No AI            & vw & 3.147 & *** & 3.104 & *** & 3.352 & *** & 2.813 & **  & no  \\
\midrule
High AI (abs)              & ew & 4.247 & *** & 4.348 & *** & 5.081 & *** & 5.075 & *** & yes \\
Moderate AI (abs)          & ew & 2.244 & **  & 2.333 & **  & 2.547 & **  & 2.887 & **  & yes \\
No AI (abs)                & ew & 1.966 & **  & 2.106 & **  & 2.271 & **  & 2.622 & **  & yes \\
S\&P~500 (abs)             & ew & 2.167 & **  & 2.522 & *** & 2.961 & *** & 3.105 & *** & no  \\
High AI / S\&P~500         & ew & 2.850 & *** & 2.947 & *** & 3.174 & *** & 3.144 & *** & yes \\
Moderate AI / S\&P~500     & ew & 1.569 & **  & 1.030 & --- & 1.632 & *   & 2.258 & **  & no  \\
No AI / S\&P~500           & ew & 1.631 & **  & 1.522 & *   & 1.481 & *   & 1.361 & --- & no  \\
High AI / No AI            & ew & 3.193 & *** & 2.985 & *** & 3.381 & *** & 2.660 & **  & no  \\
\bottomrule
\end{tabular}
\caption*{\footnotesize \textit{Notes.} Full-sample GSADF statistic and significance verdict at lag orders $k\in\{0,1,2,3\}$ for the sixteen headline series. vw and ew denote value- and equal-weighting respectively. Significance indicators $^{*}, ^{**}, ^{***}$ denote rejection of the unit-root-with-drift null at the 90\,\%, 95\,\% and 99\,\% wild-multiplier bootstrap critical values; ``---'' indicates failure to reject at 90\,\%. The ``all stable'' column reports whether the significance bin is identical at every $k$. Bootstrap critical values are recomputed at each lag, so the GSADF statistic at lag $k$ is always evaluated against the null calibrated at the same lag. Sample 2010-01 to 2026-04; $T = \num{195}$; \texttt{swindow0}~$=27$; $T_b=50$; $n_{\mathrm{boot}} = \num{2000}$.}
\end{table}


% ---------------------------------------------------------------------
\section{Wild-bootstrap multiplier}
\label{sec:rob-wild}

Inference in the PSY framework depends on bootstrap critical values for the recursive supremum statistic. \citet{harveyEtAl2016} motivate a wild-multiplier bootstrap in which the residual path is perturbed by an i.i.d. mean-zero, unit-variance multiplier sequence. The implementation used in \cref{ch:results} draws an independent $N(0,1)$ multiplier for each observation and bootstrap replication. The implementation shipped with \texttt{psymonitor}~0.0.3 instead draws one $N(0,1)$ scalar per call and recycles it across the multiplier matrix.\footnote{The relevant code is the same in the CRAN build of \texttt{psymonitor}~0.0.3 and in the package's GitHub source, and the issue was reported in the package's tracker in December~2020 without a subsequent patch.} Because the ADF $t$-statistic, and hence the BSADF and GSADF statistics, are invariant to a uniform rescaling of the regression errors, the recycled scalar is absorbed by the statistic. The package routine therefore retains residual-resampling randomness across replications, but it does not implement the observation-specific wild perturbation required by \citet{harveyEtAl2016}.

The comparison in \cref{tab:rob-multiplier} re-estimates all twenty-eight series under both multiplier procedures, holding the data, recursive specification, lag order ($k=0$), control window ($T_b=50$), bootstrap count ($n_{\mathrm{boot}}=\num{2000}$), and seed fixed. The GSADF statistic is identical across the two procedures because it is a statistic of the observed series. Only the bootstrap critical values, and therefore the significance bins, can change.

The result is that the corrected multiplier is conservative relative to the package multiplier in the cases where the classification changes. Twenty of the twenty-eight series have identical significance bins under the two procedures. The remaining eight shift downward by one bin under the corrected multiplier: five move from $^{***}$ to $^{**}$ and three move from $^{**}$ to $^{*}$. No series loses significance entirely. This pattern is consistent with the mechanism above: observation-specific multipliers allow the bootstrap innovation path to vary locally rather than being scaled by a single common draw, so the corrected procedure can generate a wider upper tail in the simulated GSADF distribution. When the corrected bootstrap raises the relevant upper critical value, a statistic that previously sat just above the package threshold may move into the adjacent lower bin.

The affected series are not the core identification contrasts. They are the Moderate~AI absolute index under both weightings, the S\&P~500 absolute benchmark under both weightings, and the Top-5~/~High~AI and Top-decile~/~High~AI ratios under both weightings. The High~AI absolute index, the High~AI~/~No~AI ratio, the High~AI~/~S\&P~500 ratio, all concentration ratios measured against No~AI, and the Mag-7~/~High~AI ratios retain the same significance classification under the corrected multiplier. The bootstrap correction therefore changes some significance labels, but it does not change the cross-sectional interpretation in \cref{ch:results}: the strongest and most thesis-relevant evidence remains located in High-AI outperformance and in mega-cap concentration, not in the secondary middle-exposure or leader-versus-High-AI contrasts.

The implication is methodological as well as substantive. The corrected $N(0,1)$ multiplier is used throughout the main results because it implements the wild bootstrap described in \citet{harveyEtAl2016}. The package multiplier would have reported eight series one significance bin higher in this application. Using the corrected multiplier therefore avoids overstating marginal evidence while leaving the main ranking of series unchanged.


% ---------------------------------------------------------------------
\section{Log-versus-level specification}
\label{sec:rob-log}

The baseline tests in \cref{ch:results} are estimated in levels, following applications such as \citet{baselePhillipsShi2025} and \citet{caspiPhillipsShi2018}. Other foundational applications, including \citet{phillipsWuYu2011} and \citet{phillipsShiYu2015a}, test logged prices or valuation ratios. This distinction matters for total-return indices. The ADF regression includes a drift term, which can absorb a linear trend in the test variable. It does not, however, absorb exponential trend in a level series. A total-return index that grows at a steady compound rate may therefore generate a positive coefficient on the lagged level even when there is no acceleration beyond ordinary compound growth. Mechanically, if a level index grows at a constant rate $g$, then $\Delta y_t \approx g y_{t-1}$, so the recursive ADF regression can read ordinary compounding as a positive coefficient on $y_{t-1}$. In logs, constant compound growth becomes approximately linear in time and is absorbed by the drift. The log specification is therefore a stricter test for acceleration relative to compound trend.

The exercise in \cref{tab:rob-loglevel} re-estimates all twenty-eight test variables after applying $y_t\mapsto\log(y_t)$ and recomputes the corrected wild-multiplier bootstrap critical values under the transformed specification. The sample, lag order ($k=0$), control window ($T_b=50$), bootstrap count ($n_{\mathrm{boot}}=\num{2000}$), and seed are held fixed relative to the baseline.

The absolute indices are the most sensitive to the transformation. All eight absolute total-return index specifications weaken in logs. The High~AI absolute index falls from the 99\,\% bin in levels to the 90\,\% bin in logs under both weightings, with the GSADF statistic declining from 4.04 to 1.44 under value-weighting and from 4.25 to 1.26 under equal-weighting. The value-weighted Moderate~AI absolute index loses significance entirely. The No~AI and S\&P~500 absolute indices also fall from the 95\,\% bin to the 90\,\% bin under both weightings. This is exactly where the level-versus-log distinction should matter most: absolute total-return indices contain strong compound growth, and the log transformation removes much of the component that a level-form ADF regression can otherwise read as explosiveness.

The relative series are more stable. The High~AI~/~No~AI ratio rejects under both specifications and both weighting schemes: it falls from $^{***}$ to $^{**}$ under value-weighting, but remains $^{***}$ under equal-weighting. The High~AI~/~S\&P~500 ratio also remains significant, with $^{***}$ under value-weighting and $^{**}$ under equal-weighting in logs. The Mag-7~/~No~AI and Mag-7~/~High~AI ratios are invariant at the 99\,\% level under both specifications and both weightings. The Top-5~/~No~AI and Top-decile~/~No~AI ratios weaken by one bin under some weightings, but continue to reject. The unstable cases are again marginal leader-versus-High-AI contrasts: the equal-weighted Top-5~/~High~AI ratio loses significance, while the equal-weighted Top-decile~/~High~AI ratio moves upward by one bin.

The log exercise therefore sharpens the interpretation of \cref{ch:results}. The level-form absolute-index rejections should not be read on their own as clean bubble evidence, because much of their strength is tied to compound total-return growth. The relative ratios are the more informative objects for the thesis question, and the key relative ratios survive the more conservative log specification. This supports the empirical strategy in \cref{subsec:bubble-test-variables}: the strongest evidence should be interpreted through High~AI outperformance relative to No~AI and the S\&P~500, and through the concentration ratios, rather than through absolute index levels alone. The date-stamped episodes in \cref{sec:results-episodes} remain baseline level-specification evidence; the log robustness exercise supports the cross-sectional rejection pattern, not a claim that every dated episode is invariant to the transformation.


% ---------------------------------------------------------------------
\section{Summary}
\label{sec:rob-summary}

The robustness exercises support the qualified reading of the bubble results. Under lag-order variation, all High-AI series reject across $k\in\{0,1,2,3\}$, the High~AI~/~No~AI ratio remains significant under both weightings, and the Mag-7~/~No~AI ratio remains the strongest concentration signal. The lag-sensitive failures are concentrated in secondary Moderate~AI~/~S\&P~500, No~AI~/~S\&P~500, and equal-weighted Top-decile~/~High~AI diagnostics. Under multiplier variation, the corrected wild bootstrap lowers the significance bin for eight of twenty-eight series relative to the package implementation, but none loses significance and none of the central High-AI or Mag-7 contrasts is reclassified. Under the log-versus-level specification, absolute total-return indices weaken substantially, while the key relative ratios continue to reject.

The main conclusion is therefore not that every PSY result is numerically invariant. Several marginal diagnostics are sensitive to lag order, multiplier implementation, or transformation. The stronger conclusion is that the thesis-relevant pattern is robust: explosive dynamics are most pronounced in High-AI relative performance and in mega-cap AI-associated concentration, while broad absolute-index evidence must be interpreted cautiously because it is partly driven by compound equity-market growth. This is the same qualified interpretation developed in \cref{sec:results-headline,sec:results-concentration}: the evidence is consistent with bubble-like dynamics around AI-exposed and AI-associated firms, but the robustness checks reinforce that the strongest claim is cross-sectional and relative, not a stand-alone verdict based on absolute price levels.
